\documentclass[journal]{IEEEtran}

\usepackage[spanish]{babel}
\usepackage[utf8]{inputenc}
\usepackage[T1]{fontenc} 
\usepackage{graphicx}
\usepackage{amsmath}
\usepackage{cite}

\begin{document}

\title{Título del Paper}
\author{Autor~Apellido,~\IEEEmembership{Member,~IEEE}}

\maketitle

\begin{abstract}
Este documento presenta un ejemplo funcional de un paper en formato IEEE.
\end{abstract}

\begin{IEEEkeywords}
IEEE, plantilla, ejemplo
\end{IEEEkeywords}

\section{Introducción}
\label{sec:intro}
Este es un ejemplo de introducción \cite{ref1}.

\section{Trabajo Relacionado}
\label{sec:related}
Trabajos previos han abordado este problema \cite{ref2}.

\section{Metodología}
\label{sec:method}
Aquí describes tu enfoque.

\section{Resultados}
\label{sec:results}
Resultados experimentales.

\section{Conclusiones}
\label{sec:conclusion}
Conclusiones finales.

\begin{thebibliography}{99}

\bibitem{ref1}
A. Autor, ``Título del artículo,'' \emph{Revista X}, vol. 1, no. 1, pp. 1--10, 2024.

\bibitem{ref2}
B. Autor, \emph{Título del libro}. Editorial, 2023.

\end{thebibliography}

\end{document}
